Considera la funció
\[
f(x)=x^3-3x^2-9x+5 .
\]

\begin{apartats}

\apartat[1,25]{1}
Estudia la monotonia de $f$ i troba'n els extrems relatius.

\begin{solucio}
$f'(x)=3x^2-6x-9=3(x-3)(x+1)$, que s'anul·la en $x=-1$ i $x=3$.\\
Signe de $f'$: positiu a $(-\infty,-1)$, negatiu a $(-1,3)$ i positiu a $(3,+\infty)$.\\
Per tant, $f$ \textbf{creix} a $(-\infty,-1)\cup(3,+\infty)$ i \textbf{decreix} a $(-1,3)$.\\
En $x=-1$ passa de créixer a decréixer: \textbf{màxim relatiu} $(-1,10)$.\\
En $x=3$ passa de decréixer a créixer: \textbf{mínim relatiu} $(3,-22)$.
\end{solucio}

\begin{nomesllarg}
\apartat{0,75}
Comprova el tipus de cada extrem amb la derivada segona.

\begin{solucio}
$f''(x)=6x-6$.\\
$f''(-1)=-12<0$: en $x=-1$ la funció és còncava, i per tant hi ha un \textbf{màxim}.\\
$f''(3)=12>0$: en $x=3$ és convexa, i per tant hi ha un \textbf{mínim}. Coincideix amb
l'estudi del signe de $f'$.
\end{solucio}
\end{nomesllarg}

\apartat[1,25]{0,75}
Demostra que la funció $g(x)=x^3+2x+7$ és creixent a tot $\mathbb{R}$.

\begin{solucio}
$g'(x)=3x^2+2$. Com que $3x^2\ge0$, es compleix $g'(x)\ge2>0$ per a tot $x$.\\
La derivada és sempre positiva, i per tant $g$ \textbf{creix a tot $\mathbb{R}$} i no té cap
extrem relatiu.
\end{solucio}

\end{apartats}
